\documentclass[conference]{IEEEtran}
\IEEEoverridecommandlockouts
\usepackage{cite}
\usepackage{amsmath,amssymb,amsfonts}
\usepackage{graphicx}
\usepackage{textcomp}
\usepackage[table]{xcolor}
\usepackage{booktabs}
\usepackage{multirow}
\usepackage{array}
\usepackage{tikz}
\usepackage{pgfplots}
\pgfplotsset{compat=1.18}
\usetikzlibrary{positioning,arrows.meta,shapes.geometric,shapes.misc,shapes.multipart,fit,backgrounds,calc}
\usepgfplotslibrary{groupplots}
\usepackage[hidelinks]{hyperref}

% ---------- palette ----------
\definecolor{ink}{HTML}{1B2433}
\definecolor{navy}{HTML}{1F3A5F}
\definecolor{teal}{HTML}{00798C}
\definecolor{plum}{HTML}{6A3F91}
\definecolor{amber}{HTML}{B7791F}
\definecolor{rose}{HTML}{D1495B}
\definecolor{litter}{HTML}{E08E0B}
\definecolor{grf}{HTML}{7B4FA0}
\definecolor{okgreen}{HTML}{1E9E5A}
\definecolor{soft}{HTML}{EEF2F7}
\definecolor{softteal}{HTML}{E6F3F5}
\definecolor{softamber}{HTML}{FBF3E4}
\definecolor{softplum}{HTML}{F2ECF7}
\definecolor{line}{HTML}{8A94A6}

\newcommand{\yes}{\tikz[baseline=-0.6ex]\fill[navy] (0,0) circle (2.1pt);}
\newcommand{\part}{\tikz[baseline=-0.6ex]{\draw[navy,line width=0.5pt] (0,0) circle (2.1pt);\fill[navy] (0,-2.1pt) arc (-90:90:2.1pt) -- cycle;}}
\newcommand{\no}{\tikz[baseline=-0.6ex]\draw[line,line width=0.5pt] (0,0) circle (2.1pt);}

% ---- ER table helpers ----
\newcommand{\pk}{\textcolor{amber!90!black}{\textbf{PK}}}
\newcommand{\fk}{\textcolor{navy}{\textbf{FK}}}
\newcommand{\uk}{\textcolor{okgreen!80!black}{\textbf{UK}}}
\newcommand{\ix}{\textcolor{line}{\textbf{IX}}}
\newcommand{\ttl}{\textcolor{rose}{\textbf{TTL}}}
\newcommand{\emb}[1]{\multicolumn{3}{@{\,}l}{\textcolor{plum}{$\triangleright$\,\textit{#1}}}\\}
\newcommand{\entity}[4]{% name, colour, width, rows
  \textcolor{white}{\textbf{\scriptsize #1}}\nodepart{second}%
  \begin{tabular}{@{\,}p{6.6mm}@{}p{#3}@{}>{\raggedleft\arraybackslash}p{19mm}@{\,}}
  #4
  \end{tabular}}

\begin{document}

\title{UrbanFix: A Closed-Loop, AI-Validated Urban Governance Platform with a Public Accountability Registry and Privacy-Preserving Research Data Access}

\author{\IEEEauthorblockN{Parth Pancholi}
\IEEEauthorblockA{\textit{School of Computing Science and Engineering}\\
\textit{VIT Bhopal University}\\
Bhopal, India\\
email address or ORCID}
\and
\IEEEauthorblockN{Gaurav Soni}
\IEEEauthorblockA{\textit{School of Computing Science and Engineering}\\
\textit{VIT Bhopal University}\\
Bhopal, India\\
email address or ORCID}
}

\maketitle

\begin{abstract}
Most digital grievance systems classify the text of a complaint and track it as a ticket. The photograph a citizen attaches is rarely checked, urgency comes from an opaque model if it is computed at all, and the records are seen only by the administration. We present UrbanFix, an urban governance platform that follows a complaint from the citizen's photograph to a verified, publicly recorded repair and makes the resulting data available for research. Citizens file without an account, confirming each report with a one-time email code. A YOLOv8 model trained for the selected category checks and localises the defect in each photo, and a weighted questionnaire produces an urgency score that can be recomputed from the answers. Staff work through a seven-stage workflow enforced on the server, which requires assignment within the ward, photo proof and a resolution receipt before a complaint is closed. All complaints and their histories appear in an open registry without personal data, and approved researchers get time-limited, logged access to an anonymised projection. We offer only categories for which a public annotated dataset exists, so every accepted report can be checked by a model. On the validation sets, mask mAP50 is 0.724 for potholes and 0.502 for litter, and box mAP50 is 0.908 for open manholes and 0.729 for graffiti, with 19 to 33 ms per image on a laptop CPU. When the model is unsure, it flags the report for a human reviewer.
\end{abstract}

\begin{IEEEkeywords}
e-governance, smart cities, civic technology, object detection, YOLOv8, human-in-the-loop AI, open data, privacy-preserving data sharing
\end{IEEEkeywords}

% =====================================================================
\section{Introduction}
In India, pothole-related road accidents rose from 3,713 in 2020 to 5,432 in 2024, an increase of 46\%, and pothole-related accidents killed 8,963 people between 2018 and 2022 \cite{morth_annual,morth_ls}. The national grievance portal CPGRAMS received between 18.7 and 26.2 lakh grievances a year from 2019 to 2024 \cite{pib22,pib25}, and in FY~2023--24 State and Union Territory governments, which handle municipal complaints, took 119 days on average to dispose of one, against 16 days for central ministries \cite{pib26} (Fig.~\ref{fig:motiv}).

\begin{figure}[t]
\centering
\begin{tikzpicture}
\begin{groupplot}[group style={group size=2 by 1, horizontal sep=9mm},
  width=0.565\columnwidth, height=3.3cm, ybar, /pgf/bar width=6pt, ymin=0,
  tick label style={font=\tiny}, label style={font=\scriptsize}, title style={font=\scriptsize, yshift=-1.5mm},
  axis line style={draw=ink!70}, major tick length=1.5pt, ymajorgrids, grid style={line width=0.2pt, draw=gray!25},
  nodes near coords, every node near coord/.append style={font=\tiny, inner sep=1pt}, enlarge x limits=0.14,
  xtick=data, x tick label style={/pgf/number format/1000 sep=}]
\nextgroupplot[title={(a) Pothole road accidents}, ymax=7200, ytick={0,2000,4000,6000}, yticklabel style={/pgf/number format/1000 sep=}]
\addplot[fill=navy, draw=none] coordinates {(2020,3713) (2021,3625) (2022,4446) (2023,5840) (2024,5432)};
\nextgroupplot[title={(b) CPGRAMS grievances (lakh)}, ymax=37, ybar=0pt, /pgf/bar width=4.5pt, legend style={font=\tiny, at={(0.5,0.98)}, anchor=north, legend columns=2, draw=none, fill=none}, nodes near coords={}]
\addplot[fill=navy, draw=none] coordinates {(2019,18.7) (2020,22.7) (2021,20.0) (2022,19.2) (2023,19.5) (2024,26.2)};
\addplot[fill=navy!35, draw=none] coordinates {(2019,8.4) (2020,10.7) (2021,10.2) (2022,8.9) (2023,6.6) (2024,3.1)};
\legend{Received, Pending b/f}
\end{groupplot}
\end{tikzpicture}
\caption{Motivation: (a) pothole-related road accidents in India \cite{morth_ls}; (b) grievances received on CPGRAMS and the backlog brought forward each year \cite{pib22,pib25}.}
\label{fig:motiv}
\end{figure}

Recent systems automate parts of this process. Some route complaint text \cite{citiwatch,smartcity}, one predicts urgency \cite{cholke}, and others classify waste images \cite{regenx} or filter unsafe photos \cite{civicmate}. None of these complaint management tools checks whether a photo shows the reported defect. Their urgency estimates cannot be traced back to an input, and their records are visible only to officials or to the person who filed. Holding a municipality to account also needs a record of who handled each step, a public outcome, and data that others can study.

UrbanFix handles a complaint as a governance record that passes through six steps in one system: report, validate, prioritise, act, prove and learn. This paper makes the following contributions:
\begin{enumerate}
\item A vision check for each category, offered only when a public, annotated street-level dataset exists for it; a YOLOv8 model \cite{yolov8} trained on that dataset localises the defect in each photo.
\item A prioritisation scheme that keeps the citizen's weighted questionnaire score and the model's evidence as separate signals and leaves the decision to a supervisor; the model cannot reject a report.
\item A workflow modelled as a guarded finite-state machine, with assignment limited to the ward, required photo proof, a rework path and an append-only history.
\item A public registry that lists every complaint with its history; the server's queries remove the citizen's identity before anything is returned.
\item A research data gateway that gives approved researchers an eight-field pseudonymised view of the data, with a granted scope, an expiry date, daily export limits and an access log.
\end{enumerate}

Section~\ref{sec:rw} covers related work and Section~\ref{sec:arch} the architecture. Section~\ref{sec:method} describes the methods, Section~\ref{sec:impl} the implementation and Section~\ref{sec:res} the results; Section~\ref{sec:conc} concludes.

% =====================================================================
\section{Related Work}\label{sec:rw}
Infra-Alert \cite{infraalert} collects geotagged photos and produces PDF reports but leaves verification and prioritisation to future work. CitiWatch \cite{citiwatch} routes complaint text with TF--IDF and a supervised classifier at 96.8\% accuracy. The framework of Cholke \textit{et al.} \cite{cholke} combines Sentence-BERT features, a multi-task urgency network and DBSCAN duplicate grouping, reporting 81\% faster response. ReGenX \cite{regenx} classifies waste photos on device (85\%), CivicMate \cite{civicmate} screens photos for unsafe content and lets citizens upvote local issues, CivicX \cite{civicx} pairs grievance analytics with plan approval and benchmarks its classifier against officers (91.4\% vs.\ 94.1\%), and Smart City Concern \cite{smartcity} uses zero-shot BART with sentiment as a proxy for urgency.

Reported accuracy for text routing lies between 90 and 97\%. As Table~\ref{tab:rw} shows, none of the seven checks the photo against the reported defect, gives an urgency that can be traced to its inputs, publishes all complaints with their histories, or shares data with researchers. Datasets with box or mask labels now exist for road damage \cite{rdd2022,roboflow}, litter \cite{taco}, open manholes \cite{roadhazards} and graffiti \cite{storm}, which makes defect-level checks feasible for these categories.

\begin{table}[t]
\caption{Capabilities of reviewed systems and UrbanFix (UF)}
\label{tab:rw}
\centering
\setlength{\tabcolsep}{3.2pt}
\renewcommand{\arraystretch}{1.05}
\footnotesize
\begin{tabular}{@{}l*{7}{c}>{\columncolor{soft}}c@{}}
\toprule
Capability & \cite{infraalert} & \cite{citiwatch} & \cite{cholke} & \cite{regenx} & \cite{civicmate} & \cite{civicx} & \cite{smartcity} & UF\\
\midrule
Account-free filing      & \no & \part & \no & \no & \no & \no & \no & \yes\\
Photo evidence           & \yes & \no & \part & \yes & \yes & \yes & \part & \yes\\
Photo--defect validation & \no & \no & \part & \no & \no & \no & \no & \yes\\
Urgency scoring          & \no & \no & \yes & \no & \no & \yes & \yes & \yes\\
Traceable urgency        & \no & \no & \no & \no & \no & \no & \no & \yes\\
Automated routing        & \part & \yes & \yes & \part & \yes & \yes & \yes & \part\\
Field proof of repair    & \yes & \no & \no & \no & \no & \no & \no & \yes\\
Resolution receipt       & \part & \no & \no & \no & \no & \no & \no & \yes\\
Public registry          & \no & \no & \no & \no & \part & \no & \no & \yes\\
Open research data       & \no & \no & \no & \no & \no & \no & \no & \yes\\
\bottomrule
\multicolumn{9}{@{}l}{\scriptsize \yes~supported \quad \part~partial \quad \no~not supported or not reported}
\end{tabular}
\end{table}

% =====================================================================
\section{System Architecture}\label{sec:arch}
Fig.~\ref{fig:arch} shows the layers. One Expo (React Native) client serves all six roles, and the server decides which set of screens a user sees. Each request passes through an Express pipeline before reaching one of seven domain services: transport checks, upload handling, JWT authentication, a per-route role guard and schema validation. The authentication step also re-checks on every call whether the account is active, whether an invitation is still pending and whether a researcher's access has expired. The intake service sends photos to a separate FastAPI inference service only after the tracking ID has been returned, so a slow or failed model call does not affect the filing. MongoDB holds seven collections. Each complaint document embeds its status history and its AI analysis.

%% ---- arch_fig ----
\begin{figure*}[t]
\centering
\begin{tikzpicture}[x=1mm, y=1mm, font=\scriptsize,
  bx/.style={draw=#1, rounded corners=1.2pt, line width=0.5pt, fill=white, align=center, inner sep=1pt, anchor=north west},
  bx/.default=navy,
  tg/.style={font=\tiny\bfseries, text=#1, anchor=south west, inner sep=0.5pt},
  tg/.default=line,
  ar/.style={-{Latex[length=1.5mm]}, line width=0.55pt, draw=#1},
  ar/.default=navy]
% ---- actors (6 x 26mm, gap 4.4)
\foreach \i/\t in {0/Citizen,1/Public observer,2/Supervisor,3/Field worker,4/Administrator,5/Researcher}{
  \node[bx=line, fill=soft, minimum width=26mm, minimum height=5mm] (a\i) at (\i*30.4,0) {\t};
  \draw[ar=line] (\i*30.4+13,-5) -- (\i*30.4+13,-9);}
% ---- presentation
\node[bx, minimum width=178mm, minimum height=13mm] (app) at (0,-9) {};
\node[anchor=north west, font=\scriptsize\bfseries, inner sep=1.5pt] at (0.5,-9.3) {Expo / React Native client (Android, iOS, web)\ \ \normalfont\tiny role-resolved shells, JWT in secure storage, deep links, camera, file export};
\foreach \i/\t in {0/{Public \& citizen: home, registry, report, track},1/{Supervisor: queue, triage, field staff},2/{Field worker: tasks, proof, attendance},3/{Administrator: analytics, people, audit},4/{Researcher: insights, query, export}}{
  \node[bx=line, fill=soft, minimum width=33mm, text width=32mm, minimum height=5mm, font=\tiny] at (2+\i*35,-15.5) {\t};}
% ---- gateway
\node[tg] at (0,-26.2) {API GATEWAY PIPELINE (Node.js 20, Express 4)};
\foreach \i/\t in {0/{Transport\\HTTPS, CORS},1/{Uploads\\Multer, $\le$3 photos},2/{Authentication\\JWT + live account checks},3/{Authorisation\\per-route role guard},4/{Validation\\Zod schemas},5/{Controllers\\7 routers, 58 endpoints},6/{Error boundary\\typed 4xx / 5xx}}{
  \node[bx, minimum width=23mm, text width=22.5mm, minimum height=8mm, font=\tiny] (p\i) at (\i*25.83,-27) {\t};}
\foreach \i [count=\j] in {0,...,5}{\draw[ar] (\i*25.83+23,-31) -- (\j*25.83,-31);}
% ---- services
\node[tg] at (0,-39.2) {DOMAIN SERVICES (bounded contexts)};
\foreach \i/\t/\c in {0/{Identity \& access\\OTP, invites, bcrypt}/navy,1/{Complaint intake\\tracking ID, urgency $u_Q$}/navy,2/{Workflow engine\\guarded FSM, Eq.~(\ref{eq:fsm})}/navy,3/{Workforce\\attendance, leave}/navy,4/{Transparency\\registry, tracker, receipts}/navy,5/{Analytics\\KPIs, heat map, override}/navy,6/{Research gateway\\projection $\pi$, quotas, log}/plum}{
  \node[bx=\c, minimum width=23mm, text width=22.5mm, minimum height=9mm, font=\tiny] (v\i) at (\i*25.83,-40) {\t};}
\draw[ar, {Latex[length=1.5mm]}-{Latex[length=1.5mm]}] (89,-22) -- node[right, font=\tiny]{REST / JSON, Bearer JWT} (89,-27);
\draw[ar] (89,-35) -- (89,-40);
% ---- intelligence tier
\node[tg=teal, anchor=north west] at (0,-71.6) {AI INFERENCE SERVICE (FastAPI, Ultralytics YOLOv8n)};
\draw[teal, dashed, rounded corners=2pt, line width=0.4pt] (0,-56) rectangle (86,-71);
\node[bx=teal, fill=softteal, minimum width=17mm, minimum height=13mm, font=\tiny] (rt) at (1,-57) {Category\\router $f_k$};
\foreach \i/\t in {0/{pothole: v8n-seg},1/{litter: v8n-seg},2/{open manhole: v8n},3/{graffiti: v8n}}{
  \node[bx=teal, minimum width=26mm, minimum height=2.9mm, font=\tiny, inner sep=0.3pt] (m\i) at (22,-57-\i*3.4) {\t};
  \draw[teal, line width=0.4pt] (18,-63.5) -- (20,-63.5) |- (22,-58.45-\i*3.4);
  \draw[teal, line width=0.4pt] (48,-58.45-\i*3.4) -| (50,-63.5);}
\node[bx=teal, fill=softteal, minimum width=33mm, text width=32mm, minimum height=13mm, font=\tiny] (pp) at (52,-57) {NMS; keep $s_j{\ge}\tau$\\$s^*$, damage share $\delta$, flag $r$\\annotated evidence\\(Eqs.~\ref{eq:smax}--\ref{eq:review})};
\draw[ar=teal] (50,-63.5) -- (52,-63.5);
% ---- data tier
\node[tg=amber, anchor=north west] at (92,-71.6) {DATA AND INTEGRATION TIER};
\draw[amber, dashed, rounded corners=2pt, line width=0.4pt] (92,-56) rectangle (178,-71);
\node[bx=amber, fill=softamber, minimum width=22mm, text width=21mm, minimum height=13mm, font=\tiny] (db) at (93,-57) {MongoDB (Mongoose)\\7 collections\\embedded history and AI evidence};
\foreach \i/\t in {0/{Object store\\photos, proof, annotations},1/{SMTP\\OTP, updates, invites},2/{PDF engine\\resolution receipts},3/{Scheduler\\expiry notices}}{
  \node[bx=amber, minimum width=12.9mm, text width=12.6mm, minimum height=13mm, font=\tiny] (d\i) at (116.5+\i*15.25,-57) {\t};}
% ---- cross-tier arrows
\draw[ar=teal, dashed] (33,-49) -- node[left, font=\tiny, text=teal, align=right]{async photos\\after ID returned} (33,-56);
\draw[ar=teal, dashed] (68.5,-56) -- (68.5,-52) -| node[pos=0.25, above, font=\tiny, text=teal]{aiAnalysis to triage} (40,-49);
\foreach \x in {108,134,160}{\draw[ar=amber] (\x,-49) -- (\x,-56);}
\end{tikzpicture}
\caption{Layered architecture of UrbanFix. Solid arrows are synchronous requests. Dashed arrows show the AI call, which is made asynchronously after the tracking ID has been returned. The security and privacy controls of Sections~\ref{sec:wf}--\ref{sec:priv} apply to all layers.}
\label{fig:arch}
\end{figure*}

% =====================================================================
\section{Methodology}\label{sec:method}

\subsection{Category selection by dataset availability}
Let $\mathcal{K}$ be the candidate complaint categories. A category $k$ is offered only if a public dataset $\mathcal{D}_k$ of street-level images with localised labels (boxes or masks), documented provenance and a usable licence exists:
\begin{equation}
\mathcal{C}=\{\,k\in\mathcal{K} : \mathcal{D}_k \text{ is public} \wedge \text{localised} \wedge \text{street-level}\,\}.
\label{eq:gate}
\end{equation}
Four of the ten categories we started with meet this condition: pothole/road damage, garbage/litter, open manhole and graffiti. For water leakage, streetlights, illegal parking, fallen trees, road signs and electric poles we found only small collections, image-level labels or data without a licence, so these were left out. The API rejects any category outside $\mathcal{C}$.

\subsection{Traceable questionnaire urgency}
Each category has five yes/no questions $q_i$ with weights $w_i\in\{1,2\}$ (2 for safety-critical). With answers $y_i\in\{0,1\}$, the urgency score and band are
\begin{equation}
U=\frac{\sum_{i=1}^{5} w_i y_i}{\sum_{i=1}^{5} w_i},\qquad
u_Q=\begin{cases}\text{High}, & U\ge 0.6\\ \text{Medium}, & 0.3\le U<0.6\\ \text{Standard}, & U<0.3.\end{cases}
\label{eq:urg}
\end{equation}
Dividing by $\sum w_i$ makes scores comparable across categories whose questions carry different weights ($\sum w_i=7$ for pothole and manhole, 6 for litter and graffiti). For a pothole, one safety-critical ``yes'' gives $U=0.29$, which is Standard, and two give $0.57$, which is Medium. The citizen sees $u_Q$ while filing, and the stored answers are enough to recompute it (Fig.~\ref{fig:urg}).

\begin{figure}[t]
\centering
\begin{tikzpicture}[x=60mm, y=1mm, font=\scriptsize]
\fill[okgreen!12] (0,-1) rectangle (0.3,13);
\fill[amber!14] (0.3,-1) rectangle (0.6,13);
\fill[rose!12] (0.6,-1) rectangle (1.0,13);
\node[font=\tiny\bfseries, text=okgreen!70!black] at (0.15,11.5) {Standard};
\node[font=\tiny\bfseries, text=amber!80!black] at (0.45,11.5) {Medium};
\node[font=\tiny\bfseries, text=rose!80!black] at (0.8,11.5) {High};
\draw[line width=0.4pt, ink!70] (0,-1) -- (1,-1);
\foreach \x in {0,0.2,0.4,0.6,0.8,1}{\draw[ink!70] (\x,-1) -- (\x,-2); \node[below, font=\tiny] at (\x,-2) {\x};}
\node[font=\scriptsize] at (0.5,-6.2) {Questionnaire score $U$ (Eq.~\ref{eq:urg})};
\node[anchor=east, font=\tiny, align=right] at (-0.01,7) {pothole, manhole\\$w=(2,2,1,1,1)$};
\node[anchor=east, font=\tiny, align=right] at (-0.01,2.5) {litter, graffiti\\$w=(2,1,1,1,1)$};
\foreach \k in {0,...,7}{\fill[navy] (\k/7,7) circle (1.3pt);}
\foreach \k in {0,...,6}{\fill[teal] (\k/6,2.5) circle (1.3pt);}
\draw[dashed, ink!60] (0.3,-1) -- (0.3,13); \draw[dashed, ink!60] (0.6,-1) -- (0.6,13);
\end{tikzpicture}
\caption{All reachable questionnaire scores for the two weight profiles, with the urgency bands. For a pothole, a single safety-critical ``yes'' scores $2/7=0.29$ and stays in Standard.}
\label{fig:urg}
\end{figure}

\subsection{Per-category vision validation}\label{sec:ai}
For a complaint of category $k$ with photos $x_1,\dots,x_n$ ($n\le3$), the model $f_k$ returns detections $\mathcal{B}_m=f_k(x_m)=\{(b_j,s_j,M_j)\}$ with boxes $b_j$, confidences $s_j$ and, for segmentation models, masks $M_j$. The service keeps detections with $s_j\ge\tau$ ($\tau=0.25$) after non-maximum suppression and computes
\begin{align}
s^{*}&=\max_{m}\max_{j} s_j, \label{eq:smax}\\
\delta&=\max_{m}\frac{\bigl|\bigcup_j M_j^{(m)}\bigr|}{H_m W_m}, \label{eq:delta}\\
r&=\mathbf{1}\left[s^{*}<\tau\right], \label{eq:review}
\end{align}
where $\delta$ is the share of the image covered by the predicted masks and $r$ is the \texttt{needsReview} flag. The suggested urgency $u_{AI}$ equals $u_Q$, raised by one band when $\delta$ is above a configurable threshold $\theta_\delta$. The model does not change the stage of a complaint. The supervisor sees $u_Q$, $u_{AI}$, $r$ and the annotated photos together and makes the decision (Fig.~\ref{fig:aipipe}). If the model misses a defect, the result is extra review work for staff; the report itself is kept.

%% ---- ai_fig ----
\begin{figure}[t]
\centering
\begin{tikzpicture}[x=1mm, y=1mm, font=\tiny,
  st/.style={draw=#1, rounded corners=1.2pt, align=center, fill=white, inner sep=1.2pt, line width=0.5pt, minimum height=7mm},
  st/.default=navy,
  ar/.style={-{Latex[length=1.4mm]}, line width=0.5pt, draw=#1}, ar/.default=navy]
\node[st, minimum width=26mm] (q) at (13,0) {Weighted questionnaire\\$U$, $u_Q$ (Eq.~\ref{eq:urg})};
\node[st, minimum width=30mm] (sv) at (66,0) {OTP verify, save, publish,\\return tracking ID};
\node[st=teal, fill=softteal, minimum width=30mm] (inf) at (66,-12) {Router $\to f_k(x_m)$:\\boxes $b_j$, masks $M_j$, $s_j$};
\node[st=teal, fill=softteal, minimum width=24mm] (ev) at (35,-12) {$s^{*}$, $\delta$\\(Eqs.~\ref{eq:smax}--\ref{eq:delta})};
\node[diamond, draw=teal, fill=softteal, aspect=2.4, inner sep=0.4pt, line width=0.5pt] (g) at (35,-24) {$s^{*}\ge\tau$\,?};
\node[st=okgreen, minimum width=22mm, minimum height=6mm] (y) at (66,-24) {$u_{AI}$ from $u_Q$ and $\delta$};
\node[st=rose, minimum width=20mm, minimum height=6mm] (n) at (8,-24) {$r=1$ (needsReview)};
\node[st=plum, fill=softplum, minimum width=82mm, minimum height=7mm] (t) at (40,-36) {Supervisor triage: photos, annotations, $u_Q$ next to $u_{AI}$, flag $r$\\$\Rightarrow$ accept or reject with remarks (human decision, Eq.~\ref{eq:fsm})};
\draw[ar] (q) -- node[above]{submit} (sv);
\draw[ar=teal, dashed] (sv) -- node[right]{async} (inf);
\draw[ar=teal] (inf) -- (ev);
\draw[ar=teal] (ev) -- (g);
\draw[ar=okgreen] (g) -- node[above]{yes} (y);
\draw[ar=rose] (g) -- node[above]{no} (n);
\draw[ar=plum] (y) -- (y |- t.north);
\draw[ar=plum] (n) -- (n |- t.north);
\draw[ar=line, dotted] (q.south) -- node[left]{$u_Q$} (13,-17) -| (n.north);
\end{tikzpicture}
\caption{AI validation and prioritisation. The supervisor sees both urgency signals; a low-confidence result sets \texttt{needsReview} and the report stays in the queue.}
\label{fig:aipipe}
\end{figure}

%% ---- er_fig ----
\begin{figure*}[t]
\centering
\begin{tikzpicture}[x=1mm, y=1mm, font=\tiny,
  e/.style={draw=#1, line width=0.6pt, inner xsep=0pt, inner ysep=1pt, anchor=north west, rectangle split, rectangle split parts=2, rectangle split part fill={#1,white}, rectangle split part align={left}, rounded corners=0.5pt},
  rl/.style={line width=0.5pt, draw=ink},
  m/.style={font=\tiny\bfseries, inner sep=0.6pt, fill=white, text=ink}]
\renewcommand{\arraystretch}{0.92}
\node[e=navy] (C) at (0,0) {\entity{\,Complaint \normalfont(complaints)}{navy}{24mm}{
\pk & \_id & ObjectId\\
\uk & trackingId & COMP-yyyymmdd-xxxxx\\
\fk\,\ix & citizenId & $\to$ User, required\\
 & title, description & string\\
\ix & category & enum $\mathcal{C}$, $|\mathcal{C}|=4$\\
\ix & ward & enum Ward 1--10\\
\ix & stage, status & enum $S$, $\varphi(S)$\\
 & urgencyLevel & enum $u_Q$\\
\fk\,\ix & assignedTo, assignedBy & $\to$ User, nullable\\
 & images, completionImages & string[$\le$3] each\\
 & closedAt, pdfReceiptUrl & date, string\\
\ix & isPublic, createdAt & bool, date\\
\emb{statusHistory[ ] \ (embedded, append-only)}
 & status, stage & enum\\
\fk & changedBy & $\to$ User, required\\
 & remarks, changedAt & string, date\\
\emb{aiAnalysis \ (embedded, 1\,:\,1)}
 & model, detected & string, bool\\
 & confidence, detections[ ] & $s^*$, $\{(b_j,s_j)\}$\\
 & damagePercent & $\delta$\\
 & suggestedUrgency, needsReview & $u_{AI}$, $r$\\
}};
\node[e=teal] (U) at (64,0) {\entity{\,User \normalfont(users)}{teal}{25mm}{
\pk & \_id & ObjectId\\
\uk & email & lowercase\\
 & name, phone & string\\
 & password & bcrypt(12), hidden\\
\ix & role, isActive & enum(5), soft delete\\
\fk & createdBy & $\to$ User\\
 & inviteToken & SHA-256, 72 h\\
 & mustSetPassword, lastLoginAt & bool, date\\
\emb{researcher \ (embedded, role = researcher)}
 & datasetScope & aggregate $|$ records\\
 & access granted / expires & $t_g$, $t_e$\\
\fk & applicationId & $\to$ ResearchApp.\\
\emb{employee \ (embedded, field $|$ supervisor)}
\uk & employeeCode & unique, sparse\\
 & ward, phone & string\\
\fk & supervisorId & $\to$ User\\
}};
\node[e=line] (O) at (70,-63) {\entity{\,Otp \normalfont(otps)}{line}{12mm}{
\ix & email & lowercase\\
 & otp & 6 digits\\
\ttl & expiresAt & 5 min\\
}};
\node[e=plum] (RA) at (128,0) {\entity{\,ResearchApplication}{plum}{23.5mm}{
\pk & \_id & ObjectId\\
\ix & email, status & pending $|$ approved\\
 & purpose, requestedDays & $\ge$100 chars, 1--180\\
\fk & reviewedBy, createdUserId & $\to$ User\\
}};
\node[e=plum] (RL) at (128,-19) {\entity{\,ResearchAccessLog}{plum}{23.5mm}{
\pk & \_id & ObjectId\\
\fk\,\ix & researcherId & $\to$ User\\
 & action, recordCount & view $|$ query $|$ export\\
\ix & filters, ip, createdAt & mixed, date $\downarrow$\\
}};
\node[e=amber] (LR) at (128,-38) {\entity{\,LeaveRequest}{amber}{23.5mm}{
\pk & \_id & ObjectId\\
\fk\,\ix & employeeId & $\to$ User\\
 & fromDate, toDate, reason & date, string\\
\fk\,\ix & decidedBy, status & $\to$ User, enum\\
}};
\node[e=amber] (AR) at (128,-57) {\entity{\,AttendanceRecord}{amber}{23.5mm}{
\pk & \_id & ObjectId\\
\fk\,\uk & employeeId, date & (emp., day) unique\\
 & checkInAt, checkOutAt, status & date, enum\\
}};
% relationships: multiplicity labels at both ends
\draw[rl] (C.east |- 0,-11.5) -- node[m, pos=0.12, above]{0..*} node[m, pos=0.88, above]{1} (U.west |- 0,-11.5);
\draw[rl] (C.east |- 0,-24) -- node[m, pos=0.12, above]{0..*} node[m, pos=0.88, above]{0..1} (U.west |- 0,-24);
\draw[rl] (C.east |- 0,-41) -- node[m, pos=0.12, above]{0..*} node[m, pos=0.88, above]{1} (U.west |- 0,-41);
\node[font=\tiny\itshape, text=line, align=center] at (57.5,-14.5) {files};
\node[font=\tiny\itshape, text=line, align=center] at (57.5,-27) {assigned};
\node[font=\tiny\itshape, text=line, align=center] at (57.5,-44) {authors event};
\draw[rl] (U.east |- 0,-5) -- node[m, pos=0.15, above]{0..1} node[m, pos=0.85, above]{0..1} (RA.west |- 0,-5);
\draw[rl] (U.east |- 0,-24) -- node[m, pos=0.15, above]{1} node[m, pos=0.85, above]{0..*} (RL.west |- 0,-24);
\draw[rl] (U.east |- 0,-43) -- node[m, pos=0.15, above]{1} node[m, pos=0.85, above]{0..*} (LR.west |- 0,-43);
\draw[rl] (U.east |- 0,-50) -- node[m, pos=0.2, above]{1} ++(6,0) |- node[m, pos=0.8, above]{0..*} (AR.west |- 0,-62);
\draw[rl, dashed] (U.south -| 79,0) -- node[m, right, font=\tiny\itshape]{email (logical)} (O.north -| 79,0);
\end{tikzpicture}
\caption{Physical data model (MongoDB, Mongoose schemas). PK, FK, UK, IX and TTL mark primary keys, references, unique, secondary and expiry indexes; rows marked $\triangleright$ begin embedded sub-documents. Multiplicities follow UML notation.}
\label{fig:er}
\end{figure*}
\subsection{Accountable workflow}\label{sec:wf}
The workflow is a deterministic finite-state machine $\mathcal{M}=(S,A,\Delta,s_0,F)$ (Fig.~\ref{fig:fsm}) with
\begin{align*}
S&=\{\textsf{sub},\textsf{acc},\textsf{asg},\textsf{wip},\textsf{prf},\textsf{cls},\textsf{rej}\},\\
A&=\{\textit{accept},\textit{reject},\textit{assign},\textit{start},\textit{proof},\textit{close},\textit{rework}\},
\end{align*}
start $s_0=\textsf{sub}$ and terminal set $F=\{\textsf{cls},\textsf{rej}\}$. A transition is applied only if its guard holds:
\begin{equation}
\Delta(s,a)\ \text{is defined} \iff s\in\mathcal{S}_a \wedge \rho(\text{actor})\in\mathcal{R}_a \wedge g_a(c),
\label{eq:fsm}
\end{equation}
where $\mathcal{S}_a$ and $\mathcal{R}_a$ are the source stages and roles allowed for action $a$, and $g_a$ requires remarks of at least 10 characters for decisions, a field worker in the complaint's ward who is active and not on approved leave for \textit{assign}, the assignee only for \textit{start} and \textit{proof}, and up to three proof photos for \textit{proof}. Each transition that is applied adds $(\text{stage},\text{status},\text{actor},\text{remarks},t)$ to the history, which is never edited, and \textit{close} emails a PDF receipt to the citizen. The public sees a coarser status given by $\varphi:S\to\{\text{Pending},\text{In Progress},\text{Resolved},\text{Rejected}\}$ with $\varphi(\textsf{sub})=\text{Pending}$, $\varphi(\textsf{cls})=\text{Resolved}$, $\varphi(\textsf{rej})=\text{Rejected}$ and In Progress otherwise.

\begin{figure}[t]
\centering
\begin{tikzpicture}[font=\scriptsize, >={Latex[length=1.5mm]},
  stg/.style={draw=#1, fill=#1!12, rounded corners=2pt, minimum width=13mm, minimum height=5.5mm, line width=0.55pt, font=\scriptsize\sffamily},
  lab/.style={font=\tiny, fill=white, inner sep=0.8pt}]
\node[stg=amber] (sub) at (0,0) {submitted};
\node[stg=plum] (acc) at (2.5,0) {accepted};
\node[stg=plum] (asg) at (5.0,0) {assigned};
\node[stg=plum] (wip) at (5.0,-1.45) {in\_progress};
\node[stg=plum] (prf) at (2.5,-1.45) {proof\_subm.};
\node[stg=okgreen, double] (cls) at (2.5,-2.75) {closed};
\node[stg=rose, double] (rej) at (0,-1.45) {rejected};
\draw[->, line width=0.5pt] (-1.15,0) -- (sub);
\draw[->] (sub) -- node[lab, above]{accept} (acc);
\draw[->] (acc) -- node[lab, above]{assign} (asg);
\draw[->] (asg) -- node[lab, right]{start} (wip);
\draw[->] (wip) -- node[lab, above]{proof} (prf);
\draw[->] (prf) -- node[lab, right]{close + receipt} (cls);
\draw[->] (sub) -- node[lab, left]{reject} (rej);
\draw[->, dashed, plum] (prf.north east) -- node[lab, sloped, above]{rework} (asg.south west);
\draw[->] (asg.north east) .. controls +(0.5,0.45) and +(-0.5,0.45) .. node[lab, above]{reassign} (asg.north west);
\node[font=\tiny, align=left, anchor=west] at (3.55,-2.75) {\textcolor{amber}{$\blacksquare$} Pending \ \textcolor{plum}{$\blacksquare$} In Progress\\ \textcolor{okgreen}{$\blacksquare$} Resolved \ \textcolor{rose}{$\blacksquare$} Rejected};
\end{tikzpicture}
\caption{Server-enforced complaint lifecycle (Eq.~\ref{eq:fsm}). Fill colour gives the public status $\varphi(s)$; double borders mark terminal stages.}
\label{fig:fsm}
\end{figure}

\subsection{Tiered disclosure and research projection}\label{sec:priv}
Different audiences receive different server-side projections of the same record (Table~\ref{tab:disc}). The registry query leaves out the citizen and staff references, and the tracker shows the filer in the history only as ``Citizen''. Researchers receive the eight-field projection
\begin{equation}
\pi(c)=\bigl(h_\kappa(c.\mathit{id}),\,k,\,u_Q,\,\varphi,\,\mathit{ward},\,\lfloor t_f\rfloor_{d},\,\lfloor t_c\rfloor_{d},\,\Delta t\bigr),
\label{eq:proj}
\end{equation}
\begin{equation}
h_\kappa(x)=\mathrm{trunc}_{48}\bigl(\mathrm{HMAC\text{-}SHA256}(\kappa,x)\bigr),\quad \Delta t=\tfrac{t_c-t_f}{86400\,\mathrm{s}},
\label{eq:hmac}
\end{equation}
where $\lfloor\cdot\rfloor_d$ truncates a timestamp to the day and $\kappa$ is a server secret. The record ID stays the same across exports, so researchers can join them, but it cannot be mapped back to the database or tracking ID without $\kappa$. Identity, free text, street address, tracking ID and photos are not included. Access is granted per application with scope $\sigma\in\{\text{aggregate},\text{records}\}$ and expiry $t_e\le t_g+180$ days, and each request $q$ is served only if
\begin{equation}
\begin{aligned}
&t_q<t_e\ \wedge\ \bigl(\text{rows}(q)\Rightarrow\sigma=\text{records}\bigr)\\
&\wedge\ n_{\mathrm{exp}}(d)\le5\ \wedge\ |q|\le5000,
\end{aligned}
\label{eq:quota}
\end{equation}
Every view, query and export is written to an access log, and CSV cells are escaped to prevent formula injection.

\begin{table}[t]
\caption{Disclosure of one complaint record by audience}
\label{tab:disc}
\centering
\setlength{\tabcolsep}{1.9pt}
\footnotesize
\begin{tabular}{@{}lcccccc@{}}
\toprule
Attribute & Admin & Staff & Registry & Tracker & Res.\,rec. & Res.\,agg.\\
\midrule
Citizen identity          & \yes & \no & \no & \no & \no & \no\\
Staff identity            & \yes & \yes & \no & \part & \no & \no\\
Record ID                 & \yes & \yes & \yes & \yes & \part & \no\\
Title, description, photos& \yes & \yes & \yes & \yes & \no & \no\\
Street location           & \yes & \yes & \yes & \yes & \part & \no\\
Category, urgency, status & \yes & \yes & \yes & \yes & \yes & \part\\
Timestamps                & \yes & \yes & \yes & \yes & \part & \part\\
History with remarks      & \yes & \yes & \part & \yes & \no & \no\\
\bottomrule
\multicolumn{7}{@{}l}{\scriptsize \part~transformed: HMAC ID, ward-level, day-level or aggregated.}
\end{tabular}
\end{table}

\subsection{Data model}
All five roles live in one \textsf{User} collection, with \textsf{researcher} and \textsf{employee} sub-documents for the role-specific fields. \textsf{Complaint} embeds \textsf{statusHistory[]} and \textsf{aiAnalysis} (Fig.~\ref{fig:er}). Staff accounts are deactivated rather than deleted so that history entries keep pointing to a valid user. OTPs expire through a TTL index, and invite tokens are stored only as SHA-256 hashes.



% =====================================================================
\begin{figure*}[t]
\centering
\newcommand{\detimg}[2]{\begin{minipage}[b]{2.75cm}\centering\includegraphics[height=2.75cm]{figs/#1}\\[1pt]{\scriptsize #2}\end{minipage}}
\newcommand{\scr}[2]{\begin{minipage}[b]{1.72cm}\centering\includegraphics[height=3.4cm]{figs/#1}\\[1pt]{\scriptsize #2}\end{minipage}}
\detimg{det_pothole_road_damage.jpg}{(a) Pothole 0.87}\hfill
\detimg{det_garbage_litter.jpg}{(b) Litter 0.88}\hfill
\detimg{det_open_manhole.jpg}{(c) Manhole 0.82}\hfill
\detimg{det_graffiti.jpg}{(d) Graffiti 0.94}\hfill
\scr{s04_report_question.jpg}{(e) Report}\hfill
\scr{s02_registry.jpg}{(f) Registry}\hfill
\scr{s08_supervisor_triage.jpg}{(g) Triage}\hfill
\scr{s18_research_profile.jpg}{(h) Research}
\caption{(a)--(d) Stored AI evidence on validation images: predicted masks or boxes with confidence. (e)--(h) App screens: weighted swipe questionnaire with live urgency, open public registry, supervisor triage, and researcher access with expiry countdown, quota and audit notice.}
\label{fig:samples}
\end{figure*}

\section{Implementation and Experimental Setup}\label{sec:impl}
The client uses Expo SDK~57 with React Native~0.86; the API uses Node.js~20, Express~4, Mongoose~8, Zod, bcrypt (cost 12), JWT (7-day) and PDFKit; the inference service uses FastAPI, PyTorch~2.14 and Ultralytics~8.4 \cite{yolov8}. Each category model is fine-tuned from COCO-pretrained YOLOv8n weights (segmentation where the dataset has masks, detection where it has boxes) at $640\times640$, batch 16, at most 60 epochs with early stopping (patience 15), default augmentation and seed 0, on an Apple M5 laptop with the MPS backend. Table~\ref{tab:data} lists the data. Training took 29 to 52 minutes per model.

Detection quality is reported with precision $P=\frac{TP}{TP+FP}$, recall $R=\frac{TP}{TP+FN}$ and average precision
\begin{equation}
\mathrm{AP}=\int_{0}^{1} p(r)\,dr,
\label{eq:ap}
\end{equation}
where $p(r)$ is the interpolated precision at recall $r$; mAP50 uses an IoU threshold of 0.5 and mAP50--95 averages over thresholds 0.50 to 0.95.

\begin{table}[t]
\caption{Datasets and models per category}
\label{tab:data}
\centering
\setlength{\tabcolsep}{2.2pt}
\footnotesize
\begin{tabular}{@{}llccc@{}}
\toprule
Category & Dataset & Labels & Train/val & Model\\
\midrule
Pothole / road damage & Roboflow \cite{roboflow} & masks & 720/60 & v8n-seg\\
Garbage / litter & TACO \cite{taco} & masks$^{\dagger}$ & 1214/214 & v8n-seg\\
Open manhole & Road Hazards \cite{roadhazards} & boxes & 1123/480 & v8n\\
Graffiti & STORM \cite{storm} & boxes & 813/209 & v8n\\
\bottomrule
\multicolumn{5}{@{}l}{\scriptsize $^{\dagger}$60 TACO classes merged into one; 1,428 of 1,500 images still available.}
\end{tabular}
\end{table}

% =====================================================================
\section{Results and Discussion}\label{sec:res}

\subsection{Validation models}
Table~\ref{tab:res} and Fig.~\ref{fig:bars} report the validation results for the best checkpoint of each model; Fig.~\ref{fig:pr} shows the precision--recall curves. The open manhole model scores highest (box mAP50 0.908), probably helped by the regular shape of manholes, although the dataset contains augmented copies of some images that may raise this figure. The pothole model reaches a mask mAP50 of 0.724 with similar precision and recall. The graffiti model has high precision (0.836) but misses faint tags (recall 0.604). Litter is the hardest category (mask mAP50 0.502): TACO's 60 object classes were merged into one and many objects are small. Given these recall values, rejecting reports on the model's say would lose real complaints, which is why Eq.~(\ref{eq:review}) only flags them. Inference takes 19 to 33~ms per photo on the laptop CPU and 6 to 20~ms on its GPU (Fig.~\ref{fig:lat}), so a GPU server is not required. Fig.~\ref{fig:samples} shows stored outputs and the matching app screens.

\begin{table}[t]
\caption{Validation results (best checkpoint)}
\label{tab:res}
\centering
\setlength{\tabcolsep}{3.4pt}
\footnotesize
\begin{tabular}{@{}lccccc@{}}
\toprule
Model (metric set) & $P$ & $R$ & mAP50 & mAP50--95 & CPU ms\\
\midrule
Pothole (mask)   & 0.725 & 0.711 & 0.724 & 0.426 & 33\\
Pothole (box)    & 0.733 & 0.682 & 0.719 & 0.447 & \\
Litter (mask)    & 0.757 & 0.434 & 0.502 & 0.317 & 24\\
Litter (box)     & 0.748 & 0.446 & 0.521 & 0.375 & \\
Open manhole (box) & 0.890 & 0.874 & \textbf{0.908} & 0.469 & 19\\
Graffiti (box)   & 0.836 & 0.604 & 0.729 & \textbf{0.526} & 20\\
\bottomrule
\end{tabular}
\end{table}

\begin{figure}[t]
\centering
\begin{tikzpicture}
\begin{axis}[
  width=\columnwidth, height=3.9cm,
  ybar, bar width=4.2pt, ymin=0, ymax=1.08,
  symbolic x coords={Pothole,Litter,Manhole,Graffiti},
  xtick=data, enlarge x limits=0.16,
  ylabel={Score}, ylabel style={font=\scriptsize}, tick label style={font=\scriptsize},
  ytick={0,0.25,0.5,0.75,1}, ymajorgrids, grid style={line width=0.2pt, draw=gray!25},
  legend style={font=\tiny, at={(0.5,1.02)}, anchor=south, legend columns=4, draw=none, /tikz/every even column/.append style={column sep=4pt}},
  nodes near coords, every node near coord/.append style={font=\tiny, rotate=90, anchor=west, inner sep=1pt},
  point meta=explicit symbolic, axis line style={draw=ink!70}, major tick length=2pt]
\addplot[fill=navy, draw=none] coordinates {(Pothole,0.725)[.73] (Litter,0.757)[.76] (Manhole,0.890)[.89] (Graffiti,0.836)[.84]};
\addplot[fill=navy!40, draw=none] coordinates {(Pothole,0.711)[.71] (Litter,0.434)[.43] (Manhole,0.874)[.87] (Graffiti,0.604)[.60]};
\addplot[fill=teal, draw=none] coordinates {(Pothole,0.724)[.72] (Litter,0.502)[.50] (Manhole,0.908)[.91] (Graffiti,0.729)[.73]};
\addplot[fill=amber, draw=none] coordinates {(Pothole,0.426)[.43] (Litter,0.317)[.32] (Manhole,0.469)[.47] (Graffiti,0.526)[.53]};
\legend{Precision, Recall, mAP50, mAP50--95}
\end{axis}
\end{tikzpicture}
\caption{Validation metrics per model (mask metrics for the segmentation models, box metrics for the detection models).}
\label{fig:bars}
\end{figure}

\begin{figure}[t]
\centering
\begin{tikzpicture}
\begin{axis}[
  width=\columnwidth, height=4.6cm, xmin=0, xmax=1, ymin=0, ymax=1.03,
  xlabel={Recall}, ylabel={Precision}, label style={font=\scriptsize}, tick label style={font=\scriptsize},
  grid=both, grid style={line width=0.2pt, draw=gray!20}, axis line style={draw=ink!70}, major tick length=2pt,
  legend style={font=\tiny, at={(0.02,0.03)}, anchor=south west, draw=none, fill=white, fill opacity=0.85, text opacity=1},
  every axis plot/.append style={line width=0.9pt, mark=none}]
\pgfplotsset{cycle list={{rose},{litter},{teal},{grf}}}
%% ---- pr_coords ----
% pothole_road_damage
\addplot coordinates {(0.000,1.000) (0.030,1.000) (0.060,1.000) (0.090,0.978) (0.120,0.978) (0.150,0.978) (0.180,0.978) (0.210,0.978) (0.240,0.961) (0.270,0.956) (0.300,0.956) (0.330,0.934) (0.360,0.927) (0.390,0.920) (0.420,0.907) (0.450,0.905) (0.480,0.890) (0.511,0.874) (0.541,0.857) (0.571,0.846) (0.601,0.834) (0.631,0.788) (0.661,0.764) (0.691,0.743) (0.721,0.714) (0.751,0.665) (0.781,0.480) (0.811,0.383) (0.841,0.281) (0.871,0.179) (0.901,0.000) (0.931,0.000) (0.961,0.000) (0.991,0.000) (0.996,0.000)};
% garbage_litter
\addplot coordinates {(0.000,1.000) (0.030,1.000) (0.060,1.000) (0.090,1.000) (0.120,1.000) (0.150,1.000) (0.180,0.985) (0.210,0.981) (0.240,0.979) (0.270,0.975) (0.300,0.956) (0.330,0.933) (0.360,0.910) (0.390,0.863) (0.420,0.804) (0.450,0.713) (0.480,0.591) (0.511,0.448) (0.541,0.324) (0.571,0.259) (0.601,0.219) (0.631,0.150) (0.661,0.082) (0.691,0.038) (0.721,0.000) (0.751,0.000) (0.781,0.000) (0.811,0.000) (0.841,0.000) (0.871,0.000) (0.901,0.000) (0.931,0.000) (0.961,0.000) (0.991,0.000) (0.996,0.000)};
% open_manhole
\addplot coordinates {(0.000,1.000) (0.030,1.000) (0.060,1.000) (0.090,1.000) (0.120,1.000) (0.150,1.000) (0.180,1.000) (0.210,1.000) (0.240,1.000) (0.270,1.000) (0.300,1.000) (0.330,1.000) (0.360,0.982) (0.390,0.976) (0.420,0.976) (0.450,0.976) (0.480,0.976) (0.511,0.976) (0.541,0.972) (0.571,0.972) (0.601,0.972) (0.631,0.972) (0.661,0.972) (0.691,0.972) (0.721,0.965) (0.751,0.957) (0.781,0.944) (0.811,0.938) (0.841,0.901) (0.871,0.896) (0.901,0.784) (0.931,0.603) (0.961,0.000) (0.991,0.000) (0.996,0.000)};
% graffiti
\addplot coordinates {(0.000,1.000) (0.030,1.000) (0.060,1.000) (0.090,1.000) (0.120,1.000) (0.150,1.000) (0.180,1.000) (0.210,0.982) (0.240,0.969) (0.270,0.969) (0.300,0.954) (0.330,0.953) (0.360,0.948) (0.390,0.947) (0.420,0.943) (0.450,0.940) (0.480,0.936) (0.511,0.935) (0.541,0.923) (0.571,0.880) (0.601,0.847) (0.631,0.774) (0.661,0.709) (0.691,0.666) (0.721,0.590) (0.751,0.548) (0.781,0.430) (0.811,0.368) (0.841,0.261) (0.871,0.177) (0.901,0.106) (0.931,0.053) (0.961,0.000) (0.991,0.000) (0.996,0.000)};
\legend{Pothole (mask) 0.724, Litter (mask) 0.502, Manhole (box) 0.908, Graffiti (box) 0.729}
\end{axis}
\end{tikzpicture}
\caption{Precision--recall curves on the validation sets; legend gives AP at IoU 0.5.}
\label{fig:pr}
\end{figure}

\begin{figure}[t]
\centering
\begin{tikzpicture}
\begin{axis}[width=\columnwidth, height=3.3cm, ybar, bar width=7pt, ymin=0, ymax=40,
  symbolic x coords={Pothole,Litter,Manhole,Graffiti}, xtick=data, enlarge x limits=0.18,
  ylabel={ms / image}, label style={font=\scriptsize}, tick label style={font=\scriptsize},
  ymajorgrids, grid style={line width=0.2pt, draw=gray!25}, axis line style={draw=ink!70}, major tick length=2pt,
  nodes near coords, every node near coord/.append style={font=\tiny, inner sep=1pt},
  legend style={font=\tiny, at={(0.98,0.97)}, anchor=north east, legend columns=2, draw=none, fill=none}]
\addplot[fill=navy, draw=none] coordinates {(Pothole,32.9) (Litter,24.2) (Manhole,18.6) (Graffiti,19.7)};
\addplot[fill=teal!70, draw=none] coordinates {(Pothole,10.7) (Litter,19.6) (Manhole,7.7) (Graffiti,6.1)};
\legend{CPU, Apple MPS GPU}
\end{axis}
\end{tikzpicture}
\caption{Mean inference time per 640\,px image including pre- and post-processing (40 images, batch 1, after warm-up).}
\label{fig:lat}
\end{figure}




\begin{figure}[t]
\centering
\begin{tikzpicture}
\begin{axis}[
  width=\columnwidth, height=3.6cm, xbar, bar width=5pt, xmin=0, xmax=96,
  y dir=reverse, symbolic y coords={Submitted,Accepted,Assigned,In progress,Proof,Closed,Rejected},
  ytick={Submitted,Accepted,Assigned,In progress,Proof,Closed,Rejected}, tick label style={font=\scriptsize}, xlabel={Complaints reaching stage}, xlabel style={font=\scriptsize},
  xmajorgrids, grid style={line width=0.2pt, draw=gray!25}, axis line style={draw=ink!70}, major tick length=2pt,
  nodes near coords, every node near coord/.append style={font=\tiny}, enlarge y limits=0.1]
\addplot[fill=navy, draw=none, bar shift=0pt] coordinates {(80,Submitted) (41,Accepted) (33,Assigned) (28,In progress) (25,Proof)};
\addplot[fill=okgreen, draw=none, bar shift=0pt] coordinates {(19,Closed)};
\addplot[fill=rose, draw=none, bar shift=0pt] coordinates {(9,Rejected)};
\end{axis}
\end{tikzpicture}
\caption{Complaints reaching each stage on the seeded demonstration dataset ($n=80$).}
\label{fig:funnel}
\end{figure}

\subsection{System validation}
We ran sixteen API tests against a live server with a separate database, and all of them passed. They cover OTP validation and expiry, rejection of categories outside $\mathcal{C}$, the absence of citizen identity in registry and tracker responses, role checks (a field worker trying to accept gets $403$), the 10-character remarks rule, a \textit{reject} after \textit{accept}, assignment across wards, a receipt request before closure, and an aggregate-only researcher asking for records ($403$). Each test targets one condition in Eqs.~(\ref{eq:fsm}) or (\ref{eq:quota}).

\subsection{Workflow on a demonstration dataset}
To test the dashboards, we seeded 80 complaints spread over twelve months and moved them through the stages (Fig.~\ref{fig:funnel}). Of these, 41 were accepted and 9 rejected at triage, and 19 were closed with a receipt, taking 8.2 to 9.2 days on average depending on the category. Because the data are synthetic, these numbers only demonstrate the pipeline and dashboards.



\subsection{Limitations}
Only four categories have suitable public data, so common complaints such as water leakage and faulty streetlights are not supported yet. The training images come from other countries, and we have not measured accuracy on Indian street photos. We chose the questionnaire weights and $\theta_\delta$ ourselves instead of fitting them to outcomes. Location is stored as a ward and free text, with no GPS and no duplicate detection. Finally, the workflow figures come from seeded data, not a field pilot.

% =====================================================================
\section{Conclusion and Future Work}\label{sec:conc}
UrbanFix handles a civic complaint from the citizen's photo to a recorded repair using open-source software and public datasets. Each photo is checked by a model trained for its category, a supervisor sets the priority with both the citizen's score and the model's evidence in view, and the server enforces each workflow step, including photo proof before closure. Records are published without personal data, and researchers get logged, pseudonymised access. Next, we plan to fine-tune the models on locally collected images, fit the questionnaire weights to closed complaints, add GPS and duplicate detection, and support regional languages and SMS reporting. We also plan a ward-level pilot with a municipal corporation to compare the model's checks with manual review.

\begin{thebibliography}{00}\footnotesize
\bibitem{morth_annual} Ministry of Road Transport and Highways, Government of India, ``Road Accidents in India,'' annual reports 2018--2022, New Delhi.
\bibitem{morth_ls} Ministry of Road Transport and Highways, reply on pothole-related road accidents 2020--2024, Lok Sabha, Feb. 12, 2026.
\bibitem{pib22} Press Information Bureau, ``Grievance cell,'' Ministry of Personnel, Public Grievances and Pensions, Mar. 16, 2022. [Online]. Available: https://pib.gov.in/PressReleasePage.aspx?PRID=1806610
\bibitem{pib25} Press Information Bureau, ``Disposal rate of grievance redressal mechanism,'' Apr. 2, 2025. [Online]. Available: https://pib.gov.in/PressReleasePage.aspx?PRID=2117810
\bibitem{pib26} Press Information Bureau, ``CPGRAMS portal,'' Feb. 5, 2026. [Online]. Available: https://pib.gov.in/PressReleasePage.aspx?PRID=2223804
\bibitem{infraalert} P.~M. Javali, S.~Athanikar, K.~Naregal, S.~Shinde, and A.~Gudnavar, ``A web-based public grievance redressal system with image and geolocation support,'' in \textit{Proc. Int. Conf. Smart \& Sustainable Technology (INCSST)}, 2025, doi: 10.1109/INCSST64791.2025.11210297.
\bibitem{citiwatch} S.~K. Kanagapriya and V.~Kesika, ``AI-driven anonymous civic intelligence for transparent and efficient municipal governance,'' in \textit{Proc. 7th Int. Conf. Intelligent Communication Technologies and Virtual Mobile Networks (ICICV)}, 2026, doi: 10.1109/ICICV68925.2026.11554648.
\bibitem{cholke} P.~Cholke, A.~S. Mishra, V.~N. Mehar, M.~R. Chikhale, H.~S. Marathe, and D.~Y. Nhalade, ``AI-driven complaint management system: An automated framework for prioritizing, tracking, and resolving public service complaints efficiently,'' in \textit{Proc. Int. Conf. Emerging Smart Computing and Informatics (ESCI)}, 2026, doi: 10.1109/ESCI68015.2026.11493166.
\bibitem{regenx} M.~Vishwanath, V.~A.~M., P.~U. Kulkarni, P.~M.~K., and V.~V. Murthy, ``AI-driven smart waste management with geofencing and civic accountability,'' in \textit{Proc. 3rd Int. Conf. Emerging Applications of Material Science and Technology (ICEAMST)}, 2025, doi: 10.1109/ICEAMST67459.2025.11336075.
\bibitem{civicmate} S.~Salunkhe, A.~Salvi, T.~Salunke, A.~Patil, A.~Saware, and R.~Jolhe, ``CivicMate: An AI-driven city grievances,'' in \textit{Proc. 6th Int. Conf. Nascent Technologies in Engineering (ICNTE)}, 2026, doi: 10.1109/ICNTE66387.2026.11437408.
\bibitem{civicx} S.~N. Katkar, V.~Parab, V.~Makadiya, M.~Pingale, and A.~Mane, ``CivicX: An AI-driven data-centric platform for urban grievance management and intelligent plan approval,'' in \textit{Proc. Int. Conf. Communication, Computing and Emerging Technologies (IC3ET)}, 2026, pp. 631--641, doi: 10.1109/IC3ET64989.2026.11467223.
\bibitem{smartcity} R.~Praveen Kumar, P.~Sam Benny, Ch.~Manideep, G.~Sushmitha, and Md.~Athfaan, ``Smart City Concern: An AI-driven multimodal civic grievance management system,'' in \textit{Proc. 6th Int. Conf. Trends in Material Science and Inventive Materials (ICTMIM)}, 2026, pp. 1894--1901, doi: 10.1109/ICTMIM68190.2026.11507401.
\bibitem{rdd2022} D.~Arya, H.~Maeda, S.~K. Ghosh, D.~Toshniwal, and Y.~Sekimoto, ``RDD2022: A multi-national image dataset for automatic road damage detection,'' arXiv:2209.08538, 2022.
\bibitem{roboflow} Roboflow Universe, ``Pothole segmentation YOLOv8 dataset, v1,'' CC BY 4.0. [Online]. Available: https://universe.roboflow.com/farzad/pothole\_segmentation\_yolov8
\bibitem{taco} P.~F. Proen\c{c}a and P.~Sim\~{o}es, ``TACO: Trash annotations in context for litter detection,'' arXiv:2003.06975, 2020.
\bibitem{roadhazards} sabidrahman, ``Pothole, cracks and open manhole (Road Hazards) dataset,'' Kaggle, 2025. [Online]. Available: https://www.kaggle.com/datasets/sabidrahman/pothole-cracks-and-openmanhole
\bibitem{storm} C.~Patrikakis, P.~Kasnesis, L.~Toumanidis, and A.~Tzitamidis, ``STORM graffiti/tagging detection dataset,'' Univ. of West Attica, Zenodo, 2019, doi: 10.5281/zenodo.3238357.
\bibitem{yolov8} G.~Jocher, A.~Chaurasia, and J.~Qiu, ``Ultralytics YOLOv8,'' 2023. [Online]. Available: https://github.com/ultralytics/ultralytics
\end{thebibliography}

\end{document}
